% ---------------------------------------------------------------
% ELL305 Term Paper -- shared preamble
% OWNER: member E. Please route all preamble changes through them.
% ---------------------------------------------------------------

% --- Page geometry: exactly as the course brief specifies ---
\usepackage[a4paper,margin=1in,bindingoffset=0.5in]{geometry}

% --- Encoding and fonts ---
\usepackage[T1]{fontenc}
\usepackage[utf8]{inputenc}
\usepackage{lmodern}
\usepackage[final]{microtype}

% --- Maths ---
\usepackage{amsmath}
\usepackage{amssymb}

% --- Tables and lists ---
\usepackage{booktabs}
\usepackage{array}
% Left-aligned paragraph column for dense comparison tables
\newcolumntype{L}[1]{>{\raggedright\arraybackslash}p{#1}}
\usepackage{multirow}
\usepackage{enumitem}
\setlist{itemsep=2pt, topsep=4pt}

% --- Graphics and diagrams ---
\usepackage{graphicx}
\usepackage{xcolor}
\usepackage{colortbl}
\usepackage{tikz}
\usetikzlibrary{arrows.meta,positioning,shapes.geometric,calc,fit,
                backgrounds,automata,decorations.pathreplacing,patterns}
\usepackage{pgfplots}
\pgfplotsset{compat=1.18}
\usepackage{forest}
\usepackage{adjustbox}
% Figure wrapper: pictures are drawn at natural size (or the scale given),
% never wider than the line.
\newenvironment{bigfig}[1][1.0]
  {\begin{adjustbox}{scale=#1, max width=\linewidth, max totalheight=0.62\textheight}}
  {\end{adjustbox}}

% --- Floats and captions ---
\usepackage{float}
\usepackage[font=small,labelfont=bf,width=0.9\textwidth]{caption}
% Captions of side-by-side figures take the width of their own minipage.
\usepackage{etoolbox}
\AtBeginEnvironment{minipage}{\captionsetup{margin=0pt}}
% Let LaTeX place figures alongside text more readily, so a figure does
% not leave a half-empty page behind it.
\renewcommand{\topfraction}{0.9}
\renewcommand{\bottomfraction}{0.8}
\renewcommand{\textfraction}{0.07}
\renewcommand{\floatpagefraction}{0.8}
\setcounter{topnumber}{3}
\setlength{\textfloatsep}{10pt plus 2pt minus 2pt}
\setlength{\floatsep}{8pt plus 2pt minus 2pt}
\setlength{\intextsep}{8pt plus 2pt minus 2pt}
\setcounter{bottomnumber}{2}
\setcounter{totalnumber}{5}

% --- Including the Turnitin report PDF as Appendix A (rule 5) ---
\usepackage{pdfpages}

% --- Units, code ---
\usepackage{siunitx}
\usepackage{listings}
\lstset{
  basicstyle=\ttfamily\footnotesize, breaklines=true, frame=single,
  numbers=left, numberstyle=\tiny\color{gray},
  keywordstyle=\color{blue!70!black},
  commentstyle=\color{green!40!black}\itshape, showstringspaces=false
}

% --- Citations: numeric, in order of first citation ---
\usepackage[numbers,square,sort&compress]{natbib}

% ---------------------------------------------------------------
% Project colour scheme -- used in every figure so all chapters
% read as one document.
% ---------------------------------------------------------------
\definecolor{mesiBlue}{RGB}{31,86,142}
\definecolor{moesiOrange}{RGB}{201,110,36}
\definecolor{coreGrey}{RGB}{90,96,104}
\definecolor{memGreen}{RGB}{46,110,76}
\definecolor{alertRed}{RGB}{176,48,48}
\definecolor{softBg}{RGB}{240,243,246}
% Accent palette: mixed hues for mind maps, decision trees and neutral
% diagram elements (MESI stays blue and MOESI stays orange wherever the
% colour carries meaning).
\definecolor{tealC}{RGB}{22,128,128}
\definecolor{plumC}{RGB}{125,62,128}
\definecolor{amberC}{RGB}{186,132,18}
\definecolor{brickC}{RGB}{176,78,56}
\definecolor{leafC}{RGB}{78,128,48}
\definecolor{skyC}{RGB}{48,112,176}
\definecolor{roseC}{RGB}{184,64,112}
\definecolor{slateC}{RGB}{84,96,120}
% Branch colours used in order by the first-level branches of a map
\colorlet{mapc1}{tealC}  \colorlet{mapc2}{plumC}  \colorlet{mapc3}{amberC}
\colorlet{mapc4}{brickC} \colorlet{mapc5}{leafC}  \colorlet{mapc6}{skyC}
\colorlet{mapc7}{roseC}  \colorlet{mapc8}{slateC} \colorlet{mapc9}{tealC}
\colorlet{mapc10}{plumC} \colorlet{mapc11}{amberC}\colorlet{mapc12}{brickC}

% ---------------------------------------------------------------
% Headings. Parts use the "top" class so a Part heading sits at the
% top of a page with its mind map and overview directly beneath it,
% rather than alone on an otherwise empty page.
% ---------------------------------------------------------------
\usepackage{titlesec}
% "straight" = no forced page break: Parts that should start a new page
% are preceded by \clearpage in main.tex; the planned Parts III--VII can
% then share pages instead of each leaving a mostly blank one.
\titleclass{\part}{straight}
\titleformat{\part}[display]
  {\centering\sffamily\bfseries\color{mesiBlue}}
  {\LARGE\MakeUppercase{\partname}~\thepart}{0.3em}
  {\rule{0.6\textwidth}{0.8pt}\\[0.35em]\huge}
  [\vspace{0.1em}\rule{0.6\textwidth}{0.8pt}]
\titlespacing*{\part}{0pt}{12pt}{10pt}

% Chapters continue on the same page instead of always starting a new one,
% which removes the half-empty page at the end of most chapters.
\titleclass{\chapter}{straight}
\titleformat{\chapter}[display]
  {\sffamily\bfseries}
  {\large\color{mesiBlue}\chaptertitlename\ \thechapter}{2pt}
  {\LARGE}
\titlespacing*{\chapter}{0pt}{16pt}{10pt}

\titleformat{\section}{\sffamily\bfseries\large}{\thesection}{0.8em}{}
\titleformat{\subsection}{\sffamily\bfseries}{\thesubsection}{0.8em}{}

% --- Callout boxes ---
\usepackage[most]{tcolorbox}
\newtcolorbox{remarkbox}[1][How this report is illustrated]{
  colback=softBg, colframe=mesiBlue, coltitle=white,
  colbacktitle=mesiBlue, fonttitle=\sffamily\bfseries,
  title={#1}, boxrule=0.6pt, arc=2pt,
  left=8pt, right=8pt, top=6pt, bottom=6pt
}
\newtcolorbox{problembox}[1][Problem Statement]{
  colback=moesiOrange!6, colframe=moesiOrange, coltitle=white,
  colbacktitle=moesiOrange, fonttitle=\sffamily\bfseries,
  title={#1}, boxrule=0.8pt, arc=2pt,
  left=8pt, right=8pt, top=6pt, bottom=6pt
}
\newtcolorbox{keybox}[1][Key idea]{
  colback=memGreen!5, colframe=memGreen, coltitle=white,
  colbacktitle=memGreen, fonttitle=\sffamily\bfseries\small,
  title={#1}, boxrule=0.6pt, arc=2pt, fontupper=\small,
  left=6pt, right=6pt, top=4pt, bottom=4pt
}
\newtcolorbox{bookbox}{
  colback=white, colframe=coreGrey!60, boxrule=0.4pt, arc=1pt,
  fontupper=\small, left=6pt, right=6pt, top=3pt, bottom=3pt,
  title={\sffamily\small Mapped to the prescribed textbook},
  coltitle=coreGrey, colbacktitle=coreGrey!10, fonttitle=\sffamily\bfseries\small
}

% ---------------------------------------------------------------
% Mind maps and decision trees (all drawn with forest).
%
% forest needs a literal \end{forest} in the source, so it must never
% be hidden inside another environment's closing hook. The \Forest
% COMMAND form is safe inside macros, which is what \sectionmap uses.
%
%  * part map / chapter map : large, grow left-to-right
%      \begin{forest} chapter map, [Chapter [Topic [point]] ] \end{forest}
%  * \sectionmap{[Section [point] [point]]}
%      a compact map placed directly under every section heading
%  * decision tree : grows top-to-bottom; label an edge with el=yes
% ---------------------------------------------------------------
\forestset{
  % Each first-level branch takes its own colour (mapc1, mapc2, ...) and
  % passes a lighter tint of it to everything below it.
  branch colour/.style={
    fill=mapc#1!78, draw=mapc#1, text=white,
    for tree={edge={thick, mapc#1!70, -{Stealth[length=6pt]}}},
    for descendants={fill=mapc#1!13, draw=mapc#1!70, text=black!85},
  },
  chapter map/.style={
    for tree={
      grow'=0, parent anchor=east, child anchor=west, anchor=west,
      align=center, font=\sffamily\footnotesize,
      draw, thick, rounded corners=3pt, inner sep=3pt,
      l sep+=8pt, s sep=2pt,
    },
    where level=0{fill=slateC!90!black, draw=slateC!90!black, text=white,
                  font=\sffamily\bfseries}{},
    where level=1{branch colour/.option=n}{},
    where level=2{font=\sffamily\scriptsize}{},
  },
  part map/.style={
    chapter map,
    where level=0{fill=mesiBlue!85!black, draw=mesiBlue!85!black}{},
  },
  section map/.style={
    for tree={
      grow'=0, parent anchor=east, child anchor=west, anchor=west,
      align=center, font=\sffamily\scriptsize,
      draw, rounded corners=2pt, inner sep=2pt,
      l sep+=6pt, s sep=1.5pt,
      edge={thick, -{Stealth[length=4pt]}},
    },
    where level=0{fill=moesiOrange!88!black, draw=moesiOrange!88!black,
                  text=white, font=\sffamily\bfseries\footnotesize}{},
    where level=1{section branch/.option=n}{},
  },
  section branch/.style={
    fill=mapc#1!15, draw=mapc#1, text=black!88,
    edge={thick, mapc#1!75, -{Stealth[length=5pt]}},
    for descendants={fill=mapc#1!8, draw=mapc#1!60,
                     edge={mapc#1!60, -{Stealth[length=4pt]}}},
  },
  decision tree/.style={
    for tree={
      grow=south, parent anchor=south, child anchor=north, anchor=north,
      align=center, font=\sffamily\footnotesize,
      draw, thick, rounded corners=3pt, inner sep=4pt,
      l sep+=12pt, s sep+=10pt,
      fill=slateC!10, draw=slateC,
      edge={thick, slateC!80, -{Stealth[length=5pt]}},
    },
    where level=1{tree branch/.option=n}{},
  },
  tree branch/.style={
    for tree={fill=mapc#1!12, draw=mapc#1,
              edge={thick, mapc#1!75, -{Stealth[length=5pt]}}},
    for descendants={if n children=0{fill=mapc#1!28}{}},
    if n children=0{fill=mapc#1!28}{},
  },
  el/.style={edge label={node[midway, fill=white, inner sep=1pt,
             font=\sffamily\tiny\bfseries, text=coreGrey]{#1}}},
}
% A compact flow map placed under each section heading.
\newcommand{\sectionmap}[1]{%
  \par\vspace{-4pt}{\centering
  \adjustbox{max width=\linewidth, max totalheight=2.2cm}{\Forest{section map, #1}}\par}\vspace{3pt}}

% --- Planned chapters that appear in the index before they are written ---
\newcommand{\plannedchapter}[1]{%
  \refstepcounter{chapter}%
  \addcontentsline{toc}{chapter}{\protect\numberline{\thechapter}#1 \textit{(in progress)}}}

% --- Handy shorthands for protocol states ---
\newcommand{\stM}{\textsf{M}}
\newcommand{\stO}{\textsf{O}}
\newcommand{\stE}{\textsf{E}}
\newcommand{\stS}{\textsf{S}}
\newcommand{\stI}{\textsf{I}}
\newcommand{\stF}{\textsf{F}}
\newcommand{\dsr}{\mathrm{DSR}}

% --- Index (rule 8): built by makeindex, which latexmk runs automatically ---
\usepackage{imakeidx}
\makeindex[intoc, columns=2, title=Index]

% --- References: hyperref then cleveref, last ---
\usepackage[hidelinks]{hyperref}
\usepackage{cleveref}

% --- Glossary (rule 8): generated by LaTeX itself, no external tool ---
\usepackage[toc, nopostdot, style=list]{glossaries}
\makenoidxglossaries
% ===============================================================
%  Glossary entries (rule 8). Printed by \printnoidxglossary in
%  main.tex; \glsaddall there makes every entry appear.
% ===============================================================
\newglossaryentry{amat}{name={AMAT}, description={Average memory access time: the
  mean number of cycles a memory access takes, counting hits and the extra time
  spent on misses (\cref{eq:amat})}}
\newglossaryentry{block}{name={cache block}, description={The unit a cache stores
  and moves, 64 bytes in this report. Also called a cache line}}
\newglossaryentry{busrd}{name={BusRd}, description={Bus read request. A cache sends
  it on a read miss to obtain a readable copy of a block}}
\newglossaryentry{busrdx}{name={BusRdX}, description={Bus read-exclusive request
  (also called read-for-ownership or read-with-intent-to-modify). A cache sends it
  on a write miss; every other copy is invalidated}}
\newglossaryentry{busupgr}{name={BusUpgr}, description={Upgrade request. A cache
  that already holds a valid shared copy sends it before writing, to invalidate
  the other copies without fetching data}}
\newglossaryentry{c2c}{name={cache-to-cache transfer}, description={A data
  response supplied by another cache instead of by memory}}
\newglossaryentry{coherence}{name={cache coherence}, description={The guarantee
  that every write to one address eventually becomes visible to all cores
  (completion) and that all cores see writes to that address in one order
  (order)}}
\newglossaryentry{directory}{name={directory protocol}, description={A coherence
  protocol in which a table at the block's home node records which caches hold
  the block, so requests go point to point instead of by broadcast}}
\newglossaryentry{dirty}{name={dirty block}, description={A block whose cached
  copy has been written and so differs from memory. It must be written back
  before memory can serve it}}
\newglossaryentry{dsr}{name={DSR}, description={Dirty Sharing Ratio: the fraction
  of a MESI run's coherence transactions that are reads served by a cache holding
  the block in \stM. Defined in \cref{ch:dsr}}}
\newglossaryentry{epoch}{name={dirty epoch}, description={The interval from the
  write that makes a block dirty to the moment its data reaches memory}}
\newglossaryentry{fwb}{name={forced writeback}, description={The memory write
  MESI performs when a cache holding a block in \stM\ supplies it to a reader and
  drops to \stS. MOESI removes it by moving to \stO}}
\newglossaryentry{invalidate}{name={write-invalidate}, description={A coherence
  policy in which a writer removes every other cached copy before writing. MSI,
  MESI and MOESI are write-invalidate protocols}}
\newglossaryentry{llc}{name={LLC}, description={Last-level cache: the cache
  shared by all cores, just above main memory}}
\newglossaryentry{mesi}{name={MESI}, description={Four-state write-invalidate
  protocol (Modified, Exclusive, Shared, Invalid). It adds \stE\ to MSI so a
  private block can be written without a bus transaction}}
\newglossaryentry{mesif}{name={MESIF}, description={Intel's variant of MESI that
  adds a Forward state naming one clean sharer as the responder}}
\newglossaryentry{moesi}{name={MOESI}, description={Five-state protocol that adds
  the Owned state to MESI, letting modified data be shared without a writeback}}
\newglossaryentry{msi}{name={MSI}, description={Three-state write-invalidate
  protocol (Modified, Shared, Invalid), the base of the family}}
\newglossaryentry{noc}{name={NoC}, description={Network on chip: the interconnect
  that carries coherence messages between cores, caches and memory}}
\newglossaryentry{openmp}{name={OpenMP}, description={A standard set of compiler
  directives and library calls for writing shared-memory parallel programs in C,
  C++ and Fortran}}
\newglossaryentry{owner}{name={owner}, description={The single cache responsible
  for a dirty block: it answers requests for the block and must write it back
  when it evicts it}}
\newglossaryentry{owned}{name={Owned state (\stO)}, description={MOESI state of a
  dirty block that other caches may also hold in \stS. Memory is stale and the
  \stO\ cache is the owner}}
\newglossaryentry{sigma}{name={survival fraction ($\sigma$)}, description={The
  fraction of a MESI run's saved writeback bytes that MOESI really saves, after
  accounting for the writebacks MOESI still performs on eviction}}
\newglossaryentry{snoop}{name={snooping protocol}, description={A coherence
  protocol in which every cache watches a shared broadcast medium (a bus) and
  reacts to the requests it sees}}
\newglossaryentry{token}{name={token coherence}, description={A coherence scheme
  in which each block has a fixed number of tokens; a cache needs one token to read
  and all tokens to write}}
\newglossaryentry{update}{name={write-update}, description={A coherence policy in
  which a writer sends the new value to every cache holding a copy, instead of
  invalidating them}}
\newglossaryentry{writeback}{name={write-back cache}, description={A cache that
  updates the level below only when a modified block is evicted}}
\newglossaryentry{writethrough}{name={write-through cache}, description={A cache
  that passes every write to the level below immediately}}


% --- Running headers ---
\usepackage{fancyhdr}
\pagestyle{fancy}
\fancyhf{}
\fancyhead[L]{\small\nouppercase{\leftmark}}
\fancyhead[R]{\small\thepage}
\renewcommand{\headrulewidth}{0.4pt}
\setlength{\headheight}{14pt}
